\documentclass{amsart}
\usepackage{amsmath, amssymb}
\title{Math 21b --- Worksheet 3: Fourier Series}
\date{}
\begin{document}
\maketitle

\textbf{Warm-up.} Recall that the inner product on $C[0,2\pi]$ is $\langle f, g \rangle = \int_0^{2\pi} f(t)g(t)\,dt$. Verify that $\cos mt$ and $\cos nt$ are orthogonal for $m \neq n$.

\bigskip

\textbf{Exercise 1.} Compute the Fourier series of:
\begin{enumerate}
    \item[(a)] $f(t) = 1$ on $[-\pi, \pi]$.
    \item[(b)] $f(t) = t^2$ on $[-\pi, \pi]$ (compute $a_0$, $a_1$, $a_2$, and the general $a_n$; note the function is even).
\end{enumerate}

\bigskip

\textbf{Exercise 2.} A square wave is defined by
\[
f(t) = \begin{cases} 1 & 0 < t < \pi \\ -1 & \pi < t < 2\pi \end{cases}
\]
extended periodically.
\begin{enumerate}
    \item[(a)] Find the Fourier series. (Hint: $f$ is odd after a shift.)
    \item[(b)] Write out the first four nonzero terms explicitly.
    \item[(c)] What does Parseval's identity say about this series?
\end{enumerate}

\bigskip

\textbf{Exercise 3 (Application --- Signal Processing).} An audio signal is modeled as $f(t) = 3\sin t + \sin 3t + 0.5\sin 5t$. If you apply a low-pass filter that removes all frequencies above $n = 2$, what signal remains? Sketch both the original and filtered signals.

\bigskip

\textbf{Reflection.} In what sense is the Fourier series a ``change of basis''? What is the basis, and what is the vector space?

\end{document}
